\documentclass{article}



\usepackage{arxiv}

\usepackage[utf8]{inputenc} % allow utf-8 input
\usepackage[T1]{fontenc}    % use 8-bit T1 fonts
\usepackage{hyperref}       % hyperlinks
\usepackage{url}            % simple URL typesetting
\usepackage{booktabs}       % professional-quality tables
\usepackage{amsmath}        % \text in math mode, align, etc.
\usepackage{amsfonts}       % blackboard math symbols
\usepackage{nicefrac}       % compact symbols for 1/2, etc.
\usepackage{microtype}      % microtypography
\usepackage{graphicx}
\usepackage{subcaption}     % subfigures
\usepackage{xcolor}         % colored TODO placeholders
\usepackage[numbers,super,sort&compress]{natbib} % Nature-style numbered superscript citations
\usepackage{doi}

% Visible placeholder for numbers/figures not yet available.
\newcommand{\TODO}[1]{\textcolor{red}{[\textbf{TODO:} #1]}}

% Figures are generated by mario_fmri (`invoke paper.*`) and copied into figures/
% (see figures/README.md). \paperfigure{name} uses figures/name.pdf, else
% figures/name.png, else a red placeholder box so the manuscript still compiles.
\graphicspath{{figures/}{./}}
\newcommand{\paperfigure}[2][\textwidth]{%
  \IfFileExists{figures/#2.pdf}%
    {\includegraphics[width=#1]{#2.pdf}}%
    {\IfFileExists{figures/#2.png}%
      {\includegraphics[width=#1]{#2.png}}%
      {\fbox{\parbox[c][0.25\textheight][c]{\dimexpr#1-2\fboxsep-2\fboxrule\relax}{\centering\TODO{missing \texttt{figures/\detokenize{#2}.pdf}}}}}}}

% Formatted for submission to Scientific Data (Nature) as a Data Descriptor.
% Section order follows the Scientific Data template: Background & Summary,
% Methods (with Code availability), Data Records, Technical Validation,
% Usage Notes, then Acknowledgements / Author contributions / Competing
% interests / Data citations / References.



% Title: 110 characters maximum including spaces (Scientific Data).
\title{CNeuromod Mario: a large-scale longitudinal fMRI dataset of naturalistic videogame play}

\date{April 6, 2026}	% Here you can change the date presented in the paper title
%\date{} 					% Or removing it

\author{ \href{https://orcid.org/0000-0002-8970-1983}{\includegraphics[scale=0.06]{orcid.pdf}\hspace{1mm}Yann Harel}~$^{1,2,\ast}$ \\
	\And
	\href{https://orcid.org/0000-0002-4391-3075}{\includegraphics[scale=0.06]{orcid.pdf}\hspace{1mm}Basile Pinsard}~$^{1,2}$ \\
    \And
	\href{https://orcid.org/0000-0001-5927-221X}{\includegraphics[scale=0.06]{orcid.pdf}\hspace{1mm}Julie A. Boyle}~$^{1,2}$ \\
    \And
	\href{https://orcid.org/0000-0002-9111-0699}{\includegraphics[scale=0.06]{orcid.pdf}\hspace{1mm}Lune Bellec}~$^{1,2}$ \\
}

\renewcommand{\shorttitle}{CNeuromod Mario data descriptor}
\renewcommand{\headeright}{Data Descriptor}
\renewcommand{\undertitle}{Data Descriptor}

%%% Add PDF metadata to help others organize their library
\hypersetup{
pdftitle={CNeuromod Mario: a large-scale longitudinal fMRI dataset of naturalistic videogame play},
pdfsubject={q-bio.NC, q-bio.QM},
pdfauthor={Yann Harel},
pdfkeywords={fMRI, Open Dataset, Naturalistic Gameplay, Super Mario Bros.},
}

\begin{document}
\maketitle

\noindent
$^{1}$Centre de Recherche de l'Institut Universitaire de Gériatrie de Montréal, Montréal, QC, Canada.\\
$^{2}$Université de Montréal, Montréal, QC, Canada.\\
$^{\ast}$Corresponding author: Yann Harel (\texttt{yharel109@gmail.com}).

\begin{abstract}
We present paired, densely sampled longitudinal functional MRI datasets in which five Courtois
NeuroMod participants played two versions of \textit{Super Mario Bros.}: \textit{mario}, the
original 8-bit game, and \textit{mariostars}, its 16-bit remake, which updates the graphics while
preserving gameplay and level design. Together they comprise 171 participant-sessions, 666 runs,
4{,}606 level repetitions, and 115.2 hours of gameplay. \textit{Mario} includes an ordered
discovery phase followed by randomized practice, whereas \textit{mariostars} independently
samples the same levels, supporting longitudinal learning analyses and within-participant
cross-task replication. Deterministic replays make every gameplay frame reconstructable from
stored controller inputs and provide frame-accurate actions, in-game events, and audiovisual
features time-locked to the blood-oxygen-level-dependent signal. We document both datasets and
validate their signal quality and annotation-driven brain maps through within- and cross-dataset
reproducibility, multivariate discriminability, and cortical interpretability. The data are
released under a CC0 waiver through the Courtois NeuroMod distribution.
\end{abstract}

% keywords can be removed
\keywords{fMRI \and Open Dataset \and Naturalistic Gameplay \and Super Mario Bros. \and Skill learning}


\section*{Background \& Summary}
\label{sec:background}

Naturalistic paradigms — movies, narratives, and video games — are increasingly used to probe
brain function under rich, continuous, ecologically valid stimulation, and to build models that
generalise across cognitive domains\cite{cneuromod2025}. Video games are a particularly
attractive naturalistic task: they couple continuous audiovisual input with goal-directed motor
action and a measurable performance signal, and they elicit learning over repeated play.
\TODO{1--2 sentences situating naturalistic / videogame fMRI and its advantages; add citations.}

The Courtois NeuroMod project\cite{cneuromod2025,boyle2025ccn} collects dense, individual-level
fMRI on a small cohort scanned repeatedly over years across many naturalistic and controlled
tasks, enabling deep, subject-specific modelling. 
\TODO{within this program we developed ways (what) of recording high-quality brain data during
action videogame play; description of shinobi.}

Within this program, the \textit{shinobi}
dataset established an annotation-based analysis pipeline for naturalistic gameplay, deriving
interpretable general-linear-model (GLM) maps from frame-accurate behavioral
events\cite{harel2025shinobi}. \TODO{1--2 sentences on what shinobi showed and why a second game
matters.}

\TODO{WHY MARIO ?}
- high quality game
- action video game, covers a space not occupied by other datasets (e.g. VGDL)
- extensive game design literature
- implicit tutorial
- easy-to-use API via gym-retro

Here we jointly describe two longitudinal tasks acquired in the same five participants.
\textit{Mario} uses the original 8-bit \textit{Super Mario Bros.} and is structured into an
ordered \emph{discovery} phase and a longer randomized \emph{practice} phase. \textit{Mariostars}
uses the 16-bit remake of the same game: its graphics are updated, while gameplay mechanics and
the 22 level designs remain intact. Its independent, randomized acquisition therefore varies
visual presentation and sampling protocol while retaining the underlying game. Together, the
tasks comprise 171 participant-sessions, 666 runs, and 4{,}606 level repetitions (115.2 hours of
play). This paired design supports longitudinal analyses of skill acquisition as well as
within-participant replication across game versions (Fig.~\ref{fig:overview}).

To establish the quality and reuse value of the dataset, we provide a technical validation of
the annotation-driven maps along three complementary axes — \emph{reproducibility},
\emph{discriminability}, and \emph{interpretability} — mirroring the shinobi analysis
pipeline\cite{harel2025shinobi}, together with behavioral learning trajectories across sessions.
These analyses are intended to demonstrate data quality and reusability, and make no claims of
new scientific findings.


\section*{Methods}
\label{sec:methods}

All processing and analysis are implemented in the open-source \texttt{mario\_fmri} Python
package and orchestrated through \texttt{invoke} tasks; every output is accompanied by a JSON
sidecar recording the git commit, parameters, and software versions. The validation pipeline
parallels the \textit{shinobi} analysis\cite{harel2025shinobi} so that results are directly
comparable across games.

\subsection*{Participants}
Five participants of the Courtois NeuroMod project (\texttt{sub-01}, \texttt{sub-02},
\texttt{sub-03}, \texttt{sub-05}, \texttt{sub-06}; aged 31--47 at recruitment in 2018)
contributed in-scanner gameplay; \texttt{sub-04} did not take part. Prior gameplay experience
varied: \texttt{sub-01} and \texttt{sub-06} had previously played \textit{Super Mario Bros.};
\texttt{sub-01} and \texttt{sub-02} were regular video-game players; \texttt{sub-03} reported no
prior video-game experience. All participants provided informed consent for the open release of
their de-identified data\cite{cneuromod2025,boyle2025ccn}. \TODO{add the approving research
ethics board / protocol number.}

\subsection*{Stimulus and apparatus}
Participants used the CNeuroMod fiber-optic MR-compatible controller\cite{harel2023controller}
to play the original \textit{Super Mario Bros.} running on a console emulator through OpenAI
\texttt{gym-retro}\cite{openai2018gymretro}, recorded at 60~Hz. Because the game is fully
deterministic, only player inputs were stored; the companion \texttt{.bk2} replay files
reconstruct every play frame-by-frame.

\subsection*{Experimental design}
We use \emph{run} for a single fMRI acquisition and \emph{repetition} for a single play of a
level (from start to completion or loss of all three lives), each corresponding to exactly one
\texttt{.bk2} replay. Twenty-two of the original game's levels were used; underwater and boss
levels were excluded because their mechanics differ markedly from the rest. The experiment
comprised two phases: a \emph{discovery} phase, in which every level was played in order with
unlimited attempts until at least one completion before advancing, and a longer \emph{practice}
phase of randomly selected levels. At the end of each run, participants completed the Flow Short
Scale (FSS-2) plus two affect items (fatigue, frustration).

\subsection*{MRI acquisition}
Imaging was performed on a 3T Siemens Prisma Fit with a 64-channel head/neck coil at the Unité
de Neuroimagerie Fonctionnelle (CRIUGM). Functional runs used an accelerated
simultaneous-multislice gradient-echo EPI sequence (HCP-style; multiband factor 4,
TR~$=1.49$~s, TE~$=37$~ms, flip angle~$52^\circ$, $2\times2\times2$~mm voxels, 60 slices,
acquisition matrix $96\times96$)\cite{xu2013multiband,glasser2016hcp}. A short reversed
phase-encode acquisition was collected for susceptibility-distortion correction. Full
acquisition details follow the CNeuroMod core protocol\cite{cneuromod2025}.

\subsection*{fMRI preprocessing}
Functional data were preprocessed with \textit{fMRIPrep} 20.2.5\cite{esteban2019fmriprep}
(slice-timing correction disabled), producing preprocessed BOLD in
\texttt{MNI152NLin2009cAsym} space together with a standard set of confound regressors.
\TODO{state the exact confound subset / nilearn load\_confounds strategy used in the GLM.}

\subsection*{Behavioral annotation}
Behavioral events were extracted from the deterministic \texttt{.bk2} replays and written as
BIDS \texttt{events.tsv} regressors. Three categories of frame-accurate events are provided:
button presses (every controller input), in-game events derived from emulator RAM (enemy kills,
deaths/hits, power-ups, coins, brick smashes, level completions), and one replay event per
repetition. Twelve gameplay conditions were modelled: six player actions (\textsc{jump},
\textsc{run/throw}, \textsc{down}, \textsc{left}, \textsc{right}, \textsc{up}) and six game
events (\textit{Kill}, \textit{Hit}, \textit{Powerup\_collected}, \textit{Coin\_collected},
\textit{Brick\_smashed}, \textit{Level\_complete}). Four continuous low-level features — frame
luminance, optical-flow magnitude, audio-envelope intensity, and per-TR button-press count —
were always included as nuisance regressors so that condition estimates are not confounded by
coarse audiovisual or motor drive. All regressors were convolved with the SPM canonical
hemodynamic response function; the design-matrix correlation structure (released as a derivative,
\texttt{invoke viz.regressor-correlations}) confirms condition estimates are not driven by
collinearity with the audiovisual/motor features.

\subsection*{Behavioral learning trajectories}
Level-matched learning trajectories were computed from the behavioral records of both datasets.
Repetitions were ordered by global play index separately within each participant, phase, and
level, then split into near-equal early, middle, and late thirds. Clear rate was computed over
decided repetitions, score as the median percentage of the maximum observed score for that level,
and time-to-clear as the median duration of cleared repetitions only. Metrics were computed per
level before levels were equally averaged within participant (Fig.~\ref{fig:overview}F--H).
For each dataset--phase--metric combination, early and late participant summaries were compared
with a one-sided paired Wilcoxon signed-rank test ($N=5$), using matched rank-biserial effect
sizes and Benjamini--Hochberg FDR correction across the nine tests. Separate released trajectory
derivatives use centred 10-repetition rolling means for completion, distance, hits, and duration.

\subsection*{Signal quality: framewise displacement and tSNR}
Two complementary signal-quality metrics were computed for every available run. Framewise
displacement (FD) was taken from the fMRIPrep confounds and summarised per run as the mean FD
and as the percentage of frames exceeding 0.5~mm, volume-weighted within participant. Temporal
signal-to-noise ratio (tSNR) was computed voxel-wise per run as the temporal mean divided by the
temporal SD, then aggregated hierarchically over runs, sessions, and participants, retaining the
across-run variance alongside the mean. Subject mean tSNR maps were projected to the inflated
fsaverage5 surface and rendered in four views (left and right lateral and medial) on a colour
scale shared by both datasets. Network summaries used the seven Yeo cortical networks and five
subcortical or cerebellar regions (Fig.~\ref{fig:qc}).

\subsection*{General linear model}
GLMs were fit with \textit{nilearn}\cite{abraham2014nilearn}. At the \emph{session level}, a
fixed-effects model combined all runs within a session (SPM HRF; 5~mm FWHM isotropic Gaussian
smoothing; AR(1) noise model; drift handled through fMRIPrep confounds), producing one beta map
and one $z$-map per condition. At the \emph{subject level}, session-level $z$-maps for each
condition were combined with a one-sample $F$-test across sessions (random effects). Statistical
maps can be thresholded with voxel-wise false-discovery-rate (FDR) control at $q<0.05$. These
session- and subject-level maps are the common input to the three map-based validation analyses
described next.

\subsection*{Reproducibility: beta-map correlations}
Map reliability was quantified as the voxel-wise Pearson correlation between full session-level
beta maps within the brain mask. For each annotation represented in both datasets, five tiers
were formed: same-participant session pairs and different-participant pairs separately within
\textit{mario} and \textit{mariostars}, and same-participant pairs across datasets. Correlations
were averaged within annotation and tier, with error bars reporting the SD of the contributing
map-pair correlations (Fig.~\ref{fig:validation}B). Directional tier comparisons used one-sided
paired Wilcoxon signed-rank tests across the 15 available annotations, with matched rank-biserial
correlations as effect sizes and Benjamini--Hochberg FDR correction across the four planned
comparisons.

\subsection*{Discriminability: multivariate decoding}
To test whether gameplay conditions evoke distinguishable spatial patterns, a linear support
vector machine\cite{pedregosa2011sklearn} was trained separately for each participant on
session-level $z$-maps to classify 30 dataset--annotation classes (the 15 annotations shared by
both datasets, split by dataset). Models used leave-one-session-out cross-validation, voxel-wise
standardisation, and ANOVA feature screening retaining the top 20\% of voxels. Fold confusion
counts were summed and row-normalised for display (Fig.~\ref{fig:validation}C). Subject-level
mean balanced accuracy was compared with the 30-class chance rate ($1/30$) using a one-sided
paired Wilcoxon signed-rank test across
the five participants; the matched-pairs rank-biserial correlation is reported as effect size.

\subsection*{Interpretability: cortical effect overlays}
Unweighted mean session contrast estimates for \textit{Kill} and \textit{Hit} were masked by the
corresponding subject-level maps using a two-sided cluster-forming threshold of $|z|>3.1$ and a
minimum extent of 10 voxels, projected to the inflated fsaverage5 surface, and displayed in four
views (Fig.~\ref{fig:validation}A). The overlay encodes support in \textit{mario},
\textit{mariostars}, or both datasets and the absolute effect magnitude; effect sign is not
shown.

\subsection*{Cross-dataset comparison with \textit{mariostars}}
A second CNeuroMod dataset, \textit{mariostars}, was acquired independently with the same five
participants and the same MRI sequence. It uses a 16-bit remake of \textit{Super Mario Bros.}
that updates the graphics while preserving the original gameplay mechanics and level designs,
and samples levels using a randomized protocol. The shared participants and game structure,
together with the altered visual presentation and independent acquisition, make
\textit{mariostars} a within-participant, cross-version replication test for annotation-driven
maps. It is shown alongside \textit{mario} throughout the Technical Validation figures
(Figs.~\ref{fig:overview}--\ref{fig:validation}); all displayed panels use empirical records or
processed derivatives from both tasks.

\subsection*{Code availability}
All analysis code is implemented in the open-source \texttt{mario\_fmri} Python package
(\url{https://github.com/courtois-neuromod/mario_fmri}, version~\TODO{tag/commit}), including the
GLM, MVPA, correlation/fingerprinting, and learning-trajectory analyses, together with the
\texttt{invoke} task runner that reproduces every figure and table. Dependencies are pinned for
reproducibility, and each output carries a JSON sidecar recording the git commit, parameters, and
software versions (Python, \textit{nilearn}, \textit{numpy}, \textit{scipy},
\textit{scikit-learn}; exact versions: \TODO{fill from provenance sidecars at release}).


\section*{Data Records}
\label{sec:datarecords}

The paired \textit{mario} and \textit{mariostars} datasets are distributed by the Courtois
NeuroMod project under a CC0 waiver and are organised following the Brain Imaging Data Structure
(BIDS)\cite{gorgolewski2016bids}
(Data Citation~1); access instructions are at \url{https://docs.cneuromod.ca}. Per-subject
gameplay totals are summarised in Table~\ref{tab:behavior}, and dataset coverage and annotation
statistics are shown in Fig.~\ref{fig:overview}. The five participants contributed 120
participant-sessions and 484 runs to \textit{mario}, compared with 51 participant-sessions and 182
runs to \textit{mariostars}. These records contain 3{,}374 and 1{,}232 level repetitions,
respectively, corresponding to 84.0 and 31.2 hours of gameplay (Fig.~\ref{fig:overview}A,B). All 22
levels are represented in both tasks, although the ordered discovery protocol produced greater
and more uneven level-wise sampling in \textit{mario} (Fig.~\ref{fig:overview}C).

\paragraph{Raw gameplay and events (\texttt{mario/}).}
Each participant folder contains, per run, a BOLD acquisition and, under \texttt{gamelogs/}, the
artefacts of every repetition: the \texttt{.bk2} replay, an \texttt{.mp4} recording, a
\texttt{\_summary.json} of game outcomes, a \texttt{\_variables.json} of per-frame emulator RAM
(including player position and button states), a \texttt{\_lowlevel.npy} of low-level
audiovisual features, an \texttt{\_events.tsv} of repetition timing, and a richer
\texttt{\_desc-annotated\_events.tsv} of the frame-accurate button-press, in-game, and replay
events described in Methods. End-of-run FSS-2 and affect responses are stored as
questionnaire events.

\paragraph{Derivatives (\texttt{mario.fmriprep/}, \texttt{processed/}).}
fMRIPrep outputs (preprocessed BOLD in \texttt{MNI152NLin2009cAsym} space, brain masks, and
confound tables) are provided under \texttt{mario.fmriprep/}. Analysis outputs generated by the
\texttt{mario\_fmri} package — session- and subject-level beta/$z$-maps, MVPA results,
beta-map correlations, and behavioral learning-trajectory tables and figures — are written to
\texttt{processed/\{analysis\}/}, each with a JSON provenance sidecar and, where applicable, a
BIDS \texttt{dataset\_description.json}.

\begin{table}[t]
  \caption{Per-subject gameplay summary across the discovery and practice phases (repetitions,
    total duration, and success rate = fraction of repetitions completed). Values reproduced from
    the dataset documentation.}
  \label{tab:behavior}
  \centering
  \small
  \begin{tabular}{lrrrrrrrr}
    \toprule
    & \multicolumn{2}{c}{Repetitions} & \multicolumn{2}{c}{Duration (h:m:s)} & \multicolumn{2}{c}{Success rate} & \multicolumn{2}{c}{Total} \\
    \cmidrule(r){2-3}\cmidrule(r){4-5}\cmidrule(r){6-7}\cmidrule(r){8-9}
    Subject & Disc. & Prac. & Disc. & Prac. & Disc. & Prac. & Reps & Dur. \\
    \midrule
    sub-01 & 230 & 567 & 03:54:27 & 09:47:11 & 0.578 & 0.781 & 797 & 13:41:38 \\
    sub-02 & 227 & 487 & 04:57:35 & 12:30:24 & 0.401 & 0.671 & 714 & 17:27:59 \\
    sub-03 & 176 & 451 & 04:49:38 & 11:57:19 & 0.432 & 0.698 & 627 & 16:46:57 \\
    sub-05 & 177 & 457 & 05:30:04 & 12:27:44 & 0.367 & 0.582 & 634 & 17:57:48 \\
    sub-06 & 134 & 468 & 04:25:41 & 13:37:40 & 0.627 & 0.857 & 602 & 18:03:22 \\
    \midrule
    Total  & 944 & 2430 & 23:37:27 & 60:20:19 & 0.481 & 0.718 & 3374 & 83:57:47 \\
    \bottomrule
  \end{tabular}
\end{table}

\begin{figure}[t]
  \centering
  \paperfigure[0.96\textwidth]{fig1_datasets_overview}
  \caption{\textbf{Acquisition, annotation, and behavioral coverage of the \textit{mario} and
    \textit{mariostars} datasets.} Data are shown for the same five participants in both tasks
    (\texttt{sub-01}, \texttt{sub-02}, \texttt{sub-03}, \texttt{sub-05}, and \texttt{sub-06}).
    (\textbf{A}, \textbf{B}) Acquisition timelines for \textit{mario} and \textit{mariostars},
    respectively. Horizontal position is cumulative run count across sessions; each block is one
    run, gaps separate sessions, and the red triangles in A mark the transition from ordered
    discovery (light red) to randomized practice (dark red). Labels at right report the number of
    sessions (s), runs (r), and gameplay hours (h) for each participant. (\textbf{C}) Mean number
    of repetitions per participant and level for \textit{mario} (red) and \textit{mariostars}
    (blue). Solid and hatched portions denote cleared and failed repetitions, respectively; error
    bars show $\pm 1$ SD of total repetitions across participants ($n=5$). (\textbf{D}) Distributions across sessions of
    annotation counts for six controller actions and four in-game events, shown separately for
    each participant on a logarithmic axis. Split violins compare \textit{mario} (brick red, left half)
    with \textit{mariostars} (deep navy, right half). (\textbf{E}) Corresponding session-wise
    distributions of median press duration for the four sustained controller actions, in seconds
    on a logarithmic axis. (\textbf{F}--\textbf{H}) Level-matched behavioral trajectories for
    cleared-repetition rate (F), median score expressed as a percentage of the level-specific
    maximum (G), and median time to clear among cleared repetitions (H). Each line represents one
    participant. Repetitions were ordered separately within each participant, phase, and level and divided
    into early, middle, and late thirds; metrics were first computed per level and then averaged
    across levels so that levels contributed equally. The first two three-bin sections of each
    panel show \textit{mario} discovery and practice, and the final, navy-shaded section shows
    \textit{mariostars} practice; the time-to-clear axis in H is reversed so that improvement is upward.
    All panels show empirical behavioral data.}
  \label{fig:overview}
\end{figure}


\section*{Technical Validation}
\label{sec:validation}

We validate the dataset by showing that the functional data are of high quality and that
annotation-driven beta maps are reproducible, discriminable, and interpretable, while the
longitudinal design captures behavioral skill acquisition. The same analyses are applied to the
independently acquired \textit{mariostars} dataset (same participants) and shown alongside
\textit{mario} in each figure as a within-cohort replication. Each subsection below reports one
analysis and carries the same name as the corresponding Methods subsection. All panels contain
empirical data; statistics are exported directly from the plotted derivatives by
\texttt{invoke paper.stats}.

\subsection*{Behavioral learning trajectories}
\label{sec:validation:traj}
The longitudinal, two-phase design makes skill acquisition directly observable, providing a
behavioral validity check on the dataset. Repetition counts per level and the discovery/practice
split are summarised in Fig.~\ref{fig:overview}A--C. After controlling level difficulty by defining
early, middle, and late thirds within each level, the mean cleared-repetition rate across
participants increased from 45.6\% to 65.5\% over \textit{mario} discovery and from 61.7\% to
74.2\% over \textit{mario} practice (Fig.~\ref{fig:overview}F). Median normalized score likewise
increased from 31.6\% to 41.4\% during discovery and from 41.8\% to 49.5\% during practice
(Fig.~\ref{fig:overview}G), while median time to clear decreased from 112.2 to 106.9~s and from
94.8 to 90.0~s, respectively (Fig.~\ref{fig:overview}H). In the independent \textit{mariostars}
practice data, early-to-late summaries were comparatively stable for cleared-repetition rate
(73.6\% to 74.9\%) and normalized score (55.6\% to 54.6\%), while time to clear decreased from
93.5 to 88.4~s. One-sided early--late tests were directionally consistent with improvement for
the principal \textit{mario} summaries, but none survived BH-FDR correction across the nine
planned tests (minimum adjusted $p=0.0938$). These five-participant results therefore document
the longitudinal behavioral content rather than establish population-level learning effects.
Beyond these summaries, per-level learning curves
(completion rate, horizontal distance, hits taken, and duration over chronological repetitions)
and spatial traces of the player's path — cleared repetitions versus deaths, localising difficult
sections — are released as derivatives (\texttt{processed/learning\_curves/} and
\texttt{processed/learning\_traces/}; \texttt{invoke traj.learning-curves}, \texttt{traj.traces}).

\subsection*{Signal quality: framewise displacement and tSNR}
\label{sec:validation:qc}
Across 464 available \textit{mario} QC runs, mean FD was 0.180~mm (median 0.177,
IQR 0.078), 67.5\% of run means were below 0.2~mm, and 3.82\% of volumes exceeded
0.5~mm. Across 180 \textit{mariostars} QC runs, the corresponding values were 0.156~mm
(median 0.160, IQR 0.080), 77.8\%, and 2.89\% (Fig.~\ref{fig:qc}A). Mean temporal
signal-to-noise ratio (tSNR) maps are shown separately for every participant and dataset
(Fig.~\ref{fig:qc}B,C), while the run, session, subject, network, and key-region summaries --
including the across-run tSNR variance -- retain the complete within-subject hierarchy
(Fig.~\ref{fig:qc}D). Mean tSNR pooled over networks and key regions was 31.4 in \textit{mario}
and 32.0 in \textit{mariostars}.

\begin{figure}[t]
  \centering
  \paperfigure[\textwidth]{fig2_signal_quality}
  \caption{\textbf{Signal quality control for \textit{mario} and \textit{mariostars}.}
    Columns correspond to the same five participants throughout. (\textbf{A}) Distributions of
    per-run mean framewise displacement (top; diamonds denote means) and percentages of volumes
    exceeding 0.5~mm (bottom; volume-weighted within participant), with datasets offset within
    each participant. The dotted top-panel reference is 0.2~mm. (\textbf{B})
    Participant-specific four-view cortical mean-tSNR maps for \textit{mario}.
    (\textbf{C}) Corresponding maps for \textit{mariostars}; both datasets use the same color
    scale. (\textbf{D})
    Within-participant tSNR summaries across runs, sessions, participants, seven cortical
    networks, and five subcortical or cerebellar regions. Small translucent points are runs,
    open points are sessions, diamonds are participant estimates, and network/region error bars
    show the square root of the across-run tSNR variance. Brick-red and deep-navy marks denote
    \textit{mario} and \textit{mariostars}, respectively.}
  \label{fig:qc}
\end{figure}

\subsection*{Reproducibility: beta-map correlations}
\label{sec:validation:repro}
Across 15 shared annotations, mean same-participant correlations exceeded different-participant
correlations in both \textit{mario} ($r=0.472$ vs. $0.221$) and \textit{mariostars}
($r=0.495$ vs. $0.247$; both one-sided paired Wilcoxon $W=120$, FDR-adjusted
$p=3.05\times10^{-5}$, matched rank-biserial $r=1.00$; Fig.~\ref{fig:validation}B).
Same-participant cross-dataset correlations averaged $r=0.461$ and also exceeded the
different-participant tier of either dataset (both $W=120$, FDR-adjusted
$p=3.05\times10^{-5}$, rank-biserial $r=1.00$), supporting replication across independently
acquired tasks while preserving participant specificity.

\subsection*{Discriminability: multivariate decoding}
\label{sec:validation:discrim}
The five participant-specific decoders achieved mean balanced accuracies of 0.814--0.831
(mean 0.826), well above the 30-class chance rate of 0.033 (one-sided paired Wilcoxon across
participants, $W=15$, $p=0.0313$, matched rank-biserial $r=1.00$;
Fig.~\ref{fig:validation}C). Strong annotation-matched diagonals occur for both datasets.
Off-diagonal structure is most visible for related motor classes, particularly directional
actions, while adjacent cross-dataset cells for the same annotation provide a direct visual
index of shared representations.

\subsection*{Interpretability: cortical effect overlays}
\label{sec:validation:interp}
The \textit{Kill} and \textit{Hit} contrasts produce spatially structured, participant-specific
cortical patterns in both datasets (Fig.~\ref{fig:validation}A). Warm and cool regions show
dataset-specific suprathreshold support, whereas purple regions identify vertices supported in
both tasks. Cross-dataset overlap, especially for \textit{Hit}, alongside individual variation
in extent, provides an anatomical complement to the quantitative map correlations without
assigning a post-hoc functional label to untested clusters.

\begin{figure}[t]
  \centering
  \paperfigure[\textwidth]{fig3_validation}
  \caption{\textbf{Technical validation} along three axes. (\textbf{A}) Interpretability:
    unweighted mean session contrast estimates for \textit{Kill} (upper row) and \textit{Hit}
    (lower row), masked by subject-level maps at a two-sided cluster-forming threshold of
    $|z|>3.1$ and a minimum extent of 10 voxels, then projected to four inflated fsaverage5 views
    (left/right lateral and medial). Red, blue, and purple encode support in \textit{mario} only,
    \textit{mariostars} only, or both, respectively; colour intensity encodes absolute contrast
    magnitude. (\textbf{B}) Reproducibility: per-annotation beta-map correlations across the
    within-subject and between-subject tiers of each dataset and the within-subject cross-dataset
    tier; error bars show $\pm 1$ SD across contributing session-map pairs. (\textbf{C})
    Discriminability: the across-participant mean row-normalised MVPA confusion matrix, accompanied
    by the five participant-specific matrices. Rows are true classes and columns predicted
    classes; each annotation's \textit{mario}/\textit{mariostars} pair is adjacent, so confusion
    within a same-annotation pair indexes cross-dataset generalisation. All panels use empirical
    derivatives.}
  \label{fig:validation}
\end{figure}

\section*{Usage Notes}
\label{sec:usage}

The dataset is intended for reuse in naturalistic-task neuroimaging, skill-learning, and
brain-encoding studies. The \texttt{mario\_fmri} package reproduces every analysis in this
descriptor through \texttt{invoke} tasks, and the companion
\href{https://github.com/courtois-neuromod/mario.tutorials}{\texttt{mario.tutorials}} notebooks
illustrate end-to-end use on a single participant, including event-based GLM analysis,
reinforcement-learning agents trained on the same gameplay, and brain-encoding models mapping
agent activations onto BOLD. Because every play is reconstructable from the deterministic
\texttt{.bk2} replays, users can derive arbitrary additional annotations beyond those provided.
The strongest reuse opportunity for replication is the independently acquired \textit{mariostars}
dataset, which shares participants and game engine.

\paragraph{Limitations.}
The cohort is small by design ($N=5$, dense rather than wide) and a single game is studied;
results emphasise within-subject, individualised modelling rather than population inference.
Gameplay actions and in-game events can co-occur, so condition regressors are partially
correlated; the discovery and practice phases are imbalanced in
duration (Table~\ref{tab:behavior}). The cross-dataset comparison with \textit{mariostars}
uses the same five participants and therefore tests replication across tasks and acquisitions,
not generalisation to an independent cohort. \TODO{add any further game-specific caveats.}


\section*{Acknowledgements}
The Courtois project on neural modelling was made possible by a generous donation from the
Courtois foundation, administered by the Fondation Institut Gériatrie Montréal at CIUSSS du
Centre-Sud-de-l'île-de-Montréal and the University of Montreal. The Courtois NeuroMod team is
based at the Centre de Recherche de l'Institut Universitaire de Gériatrie de Montréal, with
several other institutions involved. See the CNeuroMod documentation for an up-to-date list of
contributors (\url{https://docs.cneuromod.ca}). \TODO{add grant numbers / personal fellowships.}

\section*{Author contributions}
\TODO{confirm.} Y.H. designed and implemented the analyses and drafted the manuscript. B.P.
developed the acquisition and preprocessing pipelines and the data-management infrastructure.
J.A.B. coordinated participant scanning and data collection. L.B. conceived and supervised the
project and secured funding. All authors reviewed and approved the manuscript.

\section*{Competing interests}
The authors declare no competing interests.

\section*{Data Citations}
\noindent
1. Courtois NeuroMod team. \textit{cneuromod.mario}. CNeuroMod / \TODO{repository and accession or
DOI} (2026). See \url{https://docs.cneuromod.ca}.


\bibliographystyle{unsrtnat}
\bibliography{references}


\end{document}
